\documentclass[10pt,a4paper]{amsart}
\usepackage[margin=19mm]{geometry}
\usepackage{amsmath,amssymb,mathtools}
\usepackage[scr=rsfs,cal=euler]{mathalfa}
\usepackage{xfrac}
\usepackage[unicode,hidelinks,bookmarks=false]{hyperref}
\newcommand{\ie}{i.e.\ }
\newcommand{\cf}{cf.\ }
\newcommand{\resp}{respectively\ }
\newcommand{\Subsection}{\subsection}
\providecommand{\nobreakdash}{\nobreak\hbox{-}}
\setlength{\emergencystretch}{2em}
\multlinegap0pt
\usepackage{bm}
\usepackage{bbm}
\usepackage{dsfont}
\usepackage{xspace}
\usepackage{cancel}
\usepackage{mathtools}
\usepackage[shortlabels]{enumitem}
\usepackage{amsthm}
\makeatletter
\@ifundefined{theorem}{\newtheorem{theorem}{Theorem}}{}
\@ifundefined{lemma}{\newtheorem{lemma}[theorem]{Lemma}}{}
\makeatother
\def\car{\mathbf{1}}
\def\dd{\mathrm{d}}
\def\haar{\lambda}
\newcommand{\RR}{\mathbb R}
\title[Duality for Delsarte's extremal problem on locally compact A]{Duality for Delsarte's extremal problem on locally compact Abelian groups}
\author{Elena E. Berdysheva, B\'alint Farkas, Marcell Ga\'al, Mita D. Ramabulana, Szil\'ard Gy. R\'ev\'esz}
\date{Source: arXiv:2603.18287v4 (CC BY 4.0)}
\begin{document}
\maketitle

\noindent\emph{Source and licence.} Duality for Delsarte's extremal problem on locally compact Abelian groups. Version arXiv:2603.18287v4 (CC BY 4.0), distributed under the Creative Commons Attribution 4.0 International licence, as recorded in the arXiv licence metadata for this exact version and shown on the abstract page https://arxiv.org/abs/2603.18287v4. Full rights notice: licenses/arxiv-2603.18287-CC-BY-4.0.txt.\par\smallskip

\noindent\textit{Editorial scope.} Counted unit: the Lemma whose source label is \texttt{l:signswap} in the article (source0.tex lines 799--809, source label \texttt{l:signswap}), delivered together with the article's own complete original proof of it (source0.tex lines 810--844, one \texttt{proof} environment of 35 lines). No mathematics is written, repaired or re-derived: statement and proof are the article's own text line for line. Required context retained uncounted: none. Every definition, display and label of those lines is retained unchanged. Omitted from this package and not redistributed: 1--798, 845--1433. Nothing inside a retained excerpt is omitted or altered: every retained body is delivered exactly as the article's lines are written, so no omission receipt of any kind accompanies this package. Citations: the retained excerpts carry no citation command at all, so no bibliography entry of the article is transcribed and no reference is redistributed. Reference adaptation: none was needed and none was made; every reference inside the delivered unit resolves inside it, so no unresolved reference or citation is produced. Printed slips kept literal: no printed word, symbol, hypothesis or number of the counted statement or of its proof was altered. Layout and class: the article's own class is \texttt{amsart} with packages including inputenc, amsfonts, amsmath, amsthm, amssymb, bbm, xcolor, enumerate, color, ulem, babel; the wrapper is the lane's pinned \texttt{amsart} editorial wrapper at 10pt on A4, which restates the theorem-like environments the retained text needs and adds the article's own macro definitions that the retained text needs (\texttt{\textbackslash RR}, \texttt{\textbackslash car}, \texttt{\textbackslash dd}, \texttt{\textbackslash haar}), with extra wrapper packages (bm, bbm, dsfont, xspace, cancel, mathtools, enumitem). The delivered numbering is the unit's own. Assets: no retained span contains a figure environment, a graphics-inclusion command, a PDF-page inclusion, a table or a TikZ picture; no image file is copied into this package, the package declares no asset dependency and no image file is redistributed with it. Licence: version arXiv:2603.18287v4 of this article is distributed under the Creative Commons Attribution 4.0 International licence, as recorded in the arXiv licence metadata for this exact version and shown on the abstract page https://arxiv.org/abs/2603.18287v4. A full rights notice is delivered at licenses/arxiv-2603.18287-CC-BY-4.0.txt. Characters: every retained excerpt is pure ASCII, so the delivered engine, XeLaTeX with its default Latin Modern text font, reports no missing character.
\begin{lemma}[Sign swap lemma]\label{l:signswap}
	 Let $S \subseteq  G$ be a symmetric Borel set of positive Haar measure  with compact closure and $0 \not\in \overline{S}$.
Then there exists a positive definite $k \in C_c(G;\RR)$ such that $k|_S \equiv -1$.
Furthermore, given any open neighborhood $V$ of $0$ with compact closure, such that with $W:=V-V$ it holds $(W+W+W) \cap S = \emptyset$, the function~$k$ can be given with the following properties.
\begin{enumerate}[(i)]
  \item If $k(x)>0$, then $x \in W$. That is, $k_{+}$ ``lives'' on $W$.
  \item If $k(x)<0$, then $x \in S+W+W$. That is, $k_{-}$ ``lives'' on $S+W+W$.
  \item $\int_G k\dd\haar =0$ and $\frac12 \| k\|_1 = \int_G k_{+}\dd\haar = \int_G k_{-}\dd\haar = \lambda(S+W)$.
  \item $k(0)=2 \lambda(S+W)/\lambda(V)$.
  \end{enumerate}
\end{lemma}

\begin{proof} Denote the ``generalized triangle function'' as $r:=r_V:=\frac{1}{\lambda(V)} \car_V \star \car_{-V}$ and its translate by $x$ as $T_xr$, that is
\[
T_xr:=r(\cdot-x).
\]
Consider any $x \in S+ W$. Then $T_xr$ vanishes outside $S+W+W$. It is obvious that $r$ vanishes outside $W$, hence we have $r \cdot T_xr\equiv 0$, (i.e., they attain non-zero values on disjoint sets, one being a subset of $W$, and the other one a subset of $S+W+W$). It is equally easy to see that
\[
r_x := 2r-T_x r - T_{-x} r \gg 0.
\]
Indeed, $r_x = r \star (2\delta_0-\delta_x-\delta_{-x})$, where $r$ is a positive definite function as a convolution square, and the measure on the right-hand side of the convolution is an (integrally) positive definite one, too.

We now define
\[
g(y):=\int_{S+W} r_x(y) \dd \lambda(x) = \int_{S+W} (2r(y)-r(y-x)-r(y+x)) \dd\lambda(x).
\]
Let now $y \in S$. Then we also have
\[
g(y) = - \int_{S+W} (r(y-x) +r(y+x)) \dd \lambda(x) = - 2 \int_{S+W} r(y-x) \dd \lambda(x),
\]
because $r \equiv 0$ on $S+W$ and because $S$ and $S+W$ are symmetric sets. Observe that the set where $r(y-\cdot) \ne 0$ is a subset of $-W + y \subseteq  -W+S = W+S$. Thus for $y \in S$
\[
g(y)=- 2 \int_{S+W} r( y-x) \dd \lambda(x) = -2 \int_{r(y-\cdot)\ne 0} r( y-x)\dd \lambda(x) = -2 \int_G r\dd\haar = -2 \lambda(V).
\]
Therefore, $k:=\frac{1}{2\lambda(V)} g$ is constant $-1$ on $S$. Furthermore, by construction $g$, and hence $k$, is also positive definite: as each $r_x \gg 0$, their sum, or integral also satisfies the defining equations of positive definiteness.

The property (i) and (ii) are guaranteed by construction. Indeed, if $y \not\in W$, then $r(y)=0$ and the defining integral of $g$ contains only negative terms; and similarly, if $y \not\in S+W+W$, then for any $x \in S+W$ we have $y \pm x \not\in W$, and $r(y \pm x) =0$, showing that the defining integral of $g$ can have only non-negative terms. Observe that this consideration also yields that there is no $y$ with the appearance in the defining integral of $g$ both positive and negative terms; if there appears a positive value of $r(y)$, then $y \in W$ and $r(y \pm x)= 0$, for all $x \in S+W$, and if there appears a negative value $-r(y-x)<0$ then $y \in S+W+W$ and $r(y)=0$. Therefore, $g_{+}(y)=\int_{S+W} 2r(y) \dd \lambda(x) = 2 \lambda(S+W) r(y)$ and (i) follows, and $g_{-}(y)=-2 \int_{S+W} r(y-x) \dd \lambda(x)$, and also (ii) is obtained.

As each $\int r_x\dd\haar =0$, also (iii) is immediate. Finally,
\[
\frac12 \|k\|_1 = \int_G k_{+} \dd\haar= \frac{1}{2\lambda(V)} \int_G g_{+}\dd\haar = \frac{1}{2\lambda(V)} \int_G 2 \lambda(S+W) r\dd\haar = \lambda(S+W) .
\]
and
\[
k(0) = \frac{1}{2\lambda(V)} g(0) = \frac{\lambda(S+W)}{\lambda(V)}.
\]
\end{proof}

\end{document}
